\\
جدول مقادیر زیر را در نظر بگیرید:

\begin{latin}
\begin{table}[H]
  \begin{center}
    \begin{tabular}{c|c c c}
      \textbf{$x$} & $0$ & $2$ & $4$ \\
      \hline
      \textbf{$f(x)$} & 9.90 & 7.94 & 23.00 \\
    \end{tabular}
  \end{center}
\end{table}
\end{latin}

\begin{enumerate}
	\item
    ابتدا به‌روش
    کمترین مربعات
    تخمینی از
    $f(x)$
    به‌فرم
    $ax^2 + bx + c$
    بزنید.
    سپس به‌کمک رابطه به‌دست‌آمده،
    مقدار
    $\int\limits^5_{-1}f(x)dx$
    را تخمین بزنید.
    
    \item
    به‌کمک روش ذوزنقه‌ای،
    مقدار
    $\int\limits^5_{-1}f(x)dx$
    را تخمین بزنید.
\end{enumerate}
\textcolor{blue}{حل
\\
\begin{enumerate}
    \item
    \begin{align*}
    \begin{bmatrix}
        3 & 0 + 2 + 4 & 0 + 4 + 16 \\
        0 + 2 + 4 & 0 + 4 + 16 & 0 + 8 + 64 \\
        0 + 4 + 16 & 0 + 8 + 64 & 0 + 16 + 256
    \end{bmatrix}
    \begin{bmatrix}
        c \\ b \\ a
    \end{bmatrix} = \begin{bmatrix}
        9.9 + 7.94 + 23 \\
        9.9 (0) + 7.94 (2) + 23 (4)) \\
        9.9 (0) + 7.94 (4) + 23 (4^2)
    \end{bmatrix}
    \end{align*}
    با دادن دستگاه بالا به ولفرام خواهیم داشت:
    \begin{align*}
        \begin{cases}
            a \approx 2.1275 \\
            b \approx -5.235 \\
            c \approx 9.9
        \end{cases}
    \end{align*}
    پس تخمین نهایی به صورت
    $2.1275 x^2 - 5.235x + 9.9$
    خواهد بود. حال با انتگرال گیری خواهیم داشت:
    \begin{align*}
        I = (x^3 \frac{2.1275}{3} - 5.235 \frac{x^2}{2} + 9.9 x) \Biggr|^{5}_{-1} = 85.935 
    \end{align*}
    \item 
    داده‌های جدول فقط در نقاط
    \(0,2,4\)
    داده شده‌اند، در حالی که برای اجرای قاعده‌ی ذوزنقه‌ای روی بازه‌ی
    \([-1,5]\)
    دست‌کم مقدار تابع در دو سر بازه نیز لازم است. بنابراین این بخش با اطلاعات موجود تعیین‌پذیر نیست؛ باید
    \(f(-1)\)
    و
    \(f(5)\)
    (و شبکه‌ی مورد نظر) مشخص شوند.
\end{enumerate}
}
